\clearpage\pagestyle{otMain}\setlength{\otBlocEspacement}{2.5mm}\fontsize{11}{15}\selectfont\corrigesChapitre{2}{Fractions et nombres premiers}
\begin{blocLivre}{cor-2-1}\begin{corrige}{1}{Des fractions égales}\demarche{Agir de la même façon sur les deux termes.}
On multiplie les deux termes de $-3/4$ par 5 : $\res{-15/20}$. Pour $42/56$, on divise les deux termes par 14 : $\res{3/4}$.
\end{corrige}\end{blocLivre}
\begin{blocLivre}{cor-2-2}\begin{corrige}{2}{Comparer deux négatifs}\demarche{Choisir un dénominateur commun positif.}
Au dénominateur commun 18 : $-5/6=-15/18$ et $-7/9=-14/18$. Comme $-15<-14$, on obtient $\res{-5/6<-7/9}$.
\end{corrige}\end{blocLivre}
\begin{blocLivre}{cor-2-3}\begin{corrige}{3}{Additionner}\demarche{Réduire au même dénominateur.}
Le dénominateur commun 12 donne $9/12-10/12=\res{-1/12}$. Le résultat négatif est cohérent puisque $3/4<5/6$.
\end{corrige}\end{blocLivre}
\begin{blocLivre}{cor-2-4}\begin{corrige}{4}{Multiplier}\demarche{Simplifier des facteurs avant de multiplier.}
Le résultat est négatif. On simplifie $21/7=3$ et $4/10=2/5$, d’où $-2\times3/5=\res{-6/5}$.
\end{corrige}\end{blocLivre}
\begin{blocLivre}{cor-2-5}\begin{corrige}{5}{Trouver un inverse}\demarche{Transformer la division en multiplication.}
L’inverse de $-5/3$ est $\res{-3/5}$. Le quotient est $\frac23\times\frac{-3}5=\res{-2/5}$. Seul le diviseur est remplacé par son inverse.
\end{corrige}\end{blocLivre}
\begin{blocLivre}{cor-2-6}\begin{corrige}{6}{Décomposer}\demarche{Écrire chaque entier comme produit.}
$180=2^2\times3^2\times5$ et $252=2^2\times3^2\times7$. Les facteurs communs $2^2\times3^2=36$ se simplifient, donc $\res{180/252=5/7}$.
\label{fig-cor-2-6}\figureGeo[corrige]{facteurs-180}
\end{corrige}\end{blocLivre}
\begin{blocLivre}{cor-2-7}\begin{corrige}{7}{Les nombres premiers}\demarche{Tester les petits diviseurs premiers.}
La liste est $\res{71,\ 73,\ 79,\ 83,\ 89,\ 97}$. Pour tester un entier inférieur à 100, il suffit d’examiner les diviseurs premiers 2, 3, 5 et 7 : un entier composé a un facteur au plus égal à sa racine carrée. $91=7\times13$.
\end{corrige}\end{blocLivre}
\begin{blocLivre}{cor-2-8}\begin{corrige}{8}{Un calcul enchaîné}\demarche{Respecter la priorité du produit.}
La multiplication est prioritaire : $\frac23\times\frac94=\frac32$. Donc $\frac56-\frac32=\frac56-\frac96=\res{-2/3}$.
\end{corrige}\end{blocLivre}
\begin{blocLivre}{cor-2-9}\begin{corrige}{9}{Diviser une somme}\demarche{Terminer la parenthèse avant d’inverser le diviseur.}
La parenthèse vaut $1/4$. On multiplie ensuite par $8/5$ : $\frac14\times\frac85=\res{2/5}$. On peut vérifier que $(2/5)\times(5/8)=1/4$.
\end{corrige}\end{blocLivre}
\begin{blocLivre}{cor-2-10}\begin{corrige}{10}{Encadrer}\demarche{Comparer les numérateurs avec un dénominateur positif.}
$-17/5=-3{,}4$, donc $\res{-4<-17/5<-3}$. Au dénominateur 30, les trois nombres valent $-102/30$, $-100/30$, $-105/30$. L’ordre croissant est $\res{-7/2<-17/5<-10/3}$.
\end{corrige}\end{blocLivre}
\begin{blocLivre}{cor-2-11}\begin{corrige}{11}{Lire une portion du reste}\demarche{Préciser à quelle quantité s’applique chaque fraction.}
Après le premier usage, il reste $(1-2/5)\times60=36$ L. On utilise ensuite $36/3=12$ L, donc il reste $\res{24\ \mathrm L}$. La seconde fraction ne porte pas sur les 60 L initiaux.
\end{corrige}\end{blocLivre}
\begin{blocLivre}{cor-2-12}\begin{corrige}{12}{Des sachets identiques}\demarche{Chercher les diviseurs communs, puis la composition.}
Un nombre de sachets doit diviser 84 et 126. Leur facteur commun maximal est $2\times3\times7=42$. Les diviseurs communs sont $\res{1,2,3,6,7,14,21,42}$. Avec 42 sachets, chacun contient $\res2$ bleues et $\res3$ vertes.
\end{corrige}\end{blocLivre}
\begin{blocLivre}{cor-2-13}\begin{corrige}{13}{Retrouver une quantité}\demarche{Diviser par la fraction connue.}
Si $B$ est le budget, $(7/12)B=175$. Donc $B=175\times12/7=\res{300}$ euros. La part restante vaut $300-175=\res{125}$ euros. Contrôle : $(7/12)\times300=175$.
\end{corrige}\end{blocLivre}
\begin{blocLivre}{cor-2-14}\begin{corrige}{14}{Une affirmation à discuter}\demarche{Tester une fraction comprise entre 0 et 1.}
$6:(1/2)=12>6$ est un contre-exemple. Pour un nombre positif, diviser par un nombre strictement compris entre 0 et 1 revient à multiplier par un inverse supérieur à 1 ; le quotient augmente. Diviser par un nombre supérieur à 1 le diminue.
\end{corrige}\end{blocLivre}
\begin{blocLivre}{cor-2-15}\begin{corrige}{15}{Trois parts successives}\demarche{Exprimer le reste final comme fraction du total.}
Après Lina, il reste $3/4$ du total. Après Sami, il reste $1/3$ de ce reste, soit $(1/3)\times(3/4)=1/4$ du total. Donc le total vaut $36\times4=\res{144}$ euros. Lina reçoit $\res{36}$ euros, Sami $(2/3)\times108=\res{72}$ euros. $36+72+36=144$.
\end{corrige}\end{blocLivre}
